
\begin{center}{\Large\textbf{국문 초록}} \noindent\rule{\textwidth}{.12mm} \end{center}
%%%%%%%%%%%%%%%%%%%%%%%%%%%%%%%%%%%%%%%%%%%%%%%%%%%%%%%%%%%%%%%%%%%%
%%% 국문 초록 입력 
% 여기에 초록을 입력할 수 있다.
초록은 언제 쓰는가? 맨 나중에 쓰는게 국룰이다. 진짜? 당연하지. 초록은 논문의 얼굴과 같다. 선행연구 \citep{smith1987tiling}에 의하면 어쩌고 저쩌고.


%%%%% 국문 초록 입력 끝
%%%%%%%%%%%%%%%%%%%%%%%%%%%%%%%%%%%%%%%%%%%%%%%%%%%%%%%%%%%%%%%%%%%%
\vskip5pt
\noindent {\bfseries 주제어:}\ \mykeywordsko % 주제어는 메크로 \mykeywordsko 에 의해 자동으로 나타날 것이다. 

\noindent\rule{\textwidth}{.12mm}
